\documentclass[11pt,a4paper]{article}

% ----- Packages -----
\usepackage[utf8]{inputenc}
\usepackage[T1]{fontenc}
\usepackage{lmodern}
\usepackage[margin=2.5cm]{geometry}
\usepackage{amsmath,amssymb,amsfonts}
\usepackage{siunitx}
\usepackage{booktabs}
\usepackage{multirow}
\usepackage{array}
\usepackage{graphicx}
\usepackage{listings}
\usepackage{xcolor}
\usepackage{hyperref}
\usepackage{caption}
\usepackage{enumitem}
\usepackage[round]{natbib}

% ----- Hyperref + colors -----
\hypersetup{
  colorlinks=true,
  linkcolor=blue!50!black,
  urlcolor=blue!60!black,
  citecolor=teal!50!black
}

% ----- MATLAB listings style -----
\definecolor{mlcomment}{rgb}{0.13,0.55,0.13}
\definecolor{mlkeyword}{rgb}{0.0,0.0,0.7}
\definecolor{mlstring}{rgb}{0.64,0.08,0.08}
\lstdefinestyle{matlab}{
  language=Matlab,
  basicstyle=\ttfamily\footnotesize,
  keywordstyle=\color{mlkeyword}\bfseries,
  commentstyle=\color{mlcomment}\itshape,
  stringstyle=\color{mlstring},
  numbers=left, numberstyle=\tiny\color{gray},
  stepnumber=1, numbersep=6pt,
  frame=single, rulecolor=\color{gray!40},
  breaklines=true, columns=fullflexible,
  showstringspaces=false, tabsize=2
}
\lstdefinestyle{pylike}{
  basicstyle=\ttfamily\footnotesize,
  keywordstyle=\color{mlkeyword}\bfseries,
  commentstyle=\color{mlcomment}\itshape,
  numbers=left, numberstyle=\tiny\color{gray},
  frame=single, rulecolor=\color{gray!40},
  breaklines=true, columns=fullflexible,
  showstringspaces=false, tabsize=2
}

% ----- Math shorthands -----
\newcommand{\gd}{\dot{\gamma}}
\newcommand{\eff}{\mathrm{eff}}
\newcommand{\Pa}{\,\mathrm{Pa}}
\newcommand{\Pas}{\,\mathrm{Pa\,s}}
\newcommand{\mms}{\,\mathrm{mm/s}}

\title{\textbf{From rheology to slicer:\\ a parameter-prediction framework for DIW printing of CMC / sunflower-oil nanoemulsion inks}}
\author{T.~M.~C.~Rodrigues \and R.~Thir\'e \\[2pt]
\small PEMM/COPPE, Universidade Federal do Rio de Janeiro}
\date{\today}

\begin{document}
\maketitle

\begin{abstract}
\noindent
Direct Ink Writing (DIW) printers built on Fused Filament Fabrication (FFF) firmware
expose a small, FFF-centric set of process parameters --- print speed, travel speed,
flow factor, layer height, and line width --- that were never designed for
non-Newtonian, viscoelastic, recoverable ink physics. This report consolidates the
rheological, modelling, and simulation outputs already produced for the CMC ink
(C15), the CMC + sunflower-oil nanoemulsion (NE) formulation (C15~Gira~5.5), and the
commercial Bozano reference into a predictive framework that fixes each slicer
parameter from a specific, identifiable physical quantity. Three deliverables are
provided: (i) a MATLAB post-processing block converting the existing
$Q(\Delta P)$ Cross-model solver into a $v_\mathrm{print}\times k_\mathrm{flow}$
lookup table; (ii) a four-step calibration-cube protocol in Materials \& Methods
prose, ready for the upcoming printing-validation campaign; and (iii) a SAOS-derived
maximum-scaffold-height analysis based on three independent yield criteria
(LVR endpoint, $G'(\omega=1\,\mathrm{rad/s})$, and a low-frequency extrapolation to
$\omega = 0.01\,\mathrm{rad/s}$). The NE inclusion is shown to be a net
\emph{printability advantage} on every dimension: lower extrusion pressure
($\approx 4\times$), lower wall shear stress (cell-safe at $963\Pa$), \emph{and}
a wider linear-viscoelastic regime ($\gamma_\mathrm{LVR}=60.9\%$) that translates
into a $\approx 3\times$ larger predicted scaffold height under the same
sub-layer-supported stacking criterion.
\end{abstract}

% ====================================================================
\section{Framework: which datum sets which slicer parameter}
\label{sec:framework}

FFF slicers expose six parameters whose semantics, in a DIW context, no longer
match their original FFF physical meaning. Table~\ref{tab:translation} maps each
slicer field to the rheological or simulated quantity that physically governs it
for the present ink system.

\begin{table}[h]
\centering
\small
\caption{Translation map: FFF slicer field $\rightarrow$ DIW physical driver
$\rightarrow$ data source already available in the project.}
\label{tab:translation}
\begin{tabular}{p{3.0cm} p{5.0cm} p{6.8cm}}
\toprule
\textbf{FFF slicer field} & \textbf{DIW physical driver} & \textbf{Primary data source} \\
\midrule
Print speed $v_\mathrm{print}$ & Volumetric balance:\newline $Q_\mathrm{nozzle}=v_\mathrm{print}\,w\,h$ & MATLAB Cross/PL solver output $Q(\Delta P)$ (\texttt{bioprinting\_algorithm\_cross\_v2.m}) \\
Flow factor $k_\mathrm{flow}$ & Die-swell + viscoelastic recovery + syringe-wall slip & Cross exponent $m$ (swell proxy) + recovery $\%$ from \texttt{Recovery\_v1.py} + cube calibration \\
Travel speed $v_\mathrm{travel}$ & Recovery time vs.\ gravitational sag time & Recovery test ($R_3/R_1$) at deposition shear; $\eta_0$ for sag \\
Layer height $h_\mathrm{layer}$ & LVR yield stress / quasi-static $G'$ & SAOS frequency sweep + amplitude sweep ($\gamma_\mathrm{LVR}$) \\
Line width $w_\mathrm{line}$ & Nozzle ID + die-swell ratio $(1+\beta)$ & Cross exponent $m$ (swell proxy) + experimental cube \\
Extrusion pressure $P_\mathrm{ext}$ & Generalised Hagen--Poiseuille with fitted Cross/PL model & Existing MATLAB solver outputs \\
\bottomrule
\end{tabular}
\end{table}

The structural difference from FFF is that here $v_\mathrm{print}$ and the flow
output are \emph{physically decoupled}: $v_\mathrm{print}$ is purely a head-motion
command, whereas the volumetric flow at the nozzle is governed by the
pressure--rheology loop. The slicer firmware still treats them as coupled through
an effective extruder ``E-steps'' calibration, which is why $k_\mathrm{flow}$ must
be measured experimentally per ink rather than assumed unity.

% ====================================================================
\section{Governing equations}
\label{sec:equations}

\subsection{Volumetric balance at the nozzle}
For a deposited rectangular filament of width $w$ and height $h$ moving with head
speed $v_\mathrm{print}$, conservation of volume at the nozzle gives
\begin{equation}
Q_\mathrm{nozzle}(\Delta P) \;=\; v_\mathrm{print}\,w\,h.
\label{eq:vol_balance}
\end{equation}
The MATLAB Cross-model solver already returns $Q(\Delta P)$ from a numerical
integration of $v(r)$ over the needle cross-section. Inverting
Eq.~\ref{eq:vol_balance} yields the maximum head speed compatible with $w \approx
\mathrm{ID}$ and a target layer height.

\subsection{Power-law analytical bound}
For an analytical sanity check on the Cross-model output, the Power-Law mean
nozzle velocity in a straight cylinder of radius $R$ and length $L$ is
\begin{equation}
\bar v_\mathrm{nozzle} \;=\; \frac{n}{3n+1}\,R\,
\left(\frac{R\,\Delta P}{2KL}\right)^{1/n}.
\label{eq:vmean_PL}
\end{equation}
With $K$ and $n$ from \texttt{FitAll-Ramp1-v4.txt} (Table~\ref{tab:rheo_params}),
Eq.~\ref{eq:vmean_PL} reproduces the Cross-model values to within $5\%$ for the
$\dot\gamma$ range that dominates the needle interior.

\subsection{Flow factor decomposition}
The slicer's $k_\mathrm{flow}$ multiplier corrects three independent deviations
from ideal FFF extrusion:
\begin{equation}
k_\mathrm{flow} \;=\; \underbrace{(1+\beta_\mathrm{swell})^{2}}_{\text{die swell}}
\cdot
\underbrace{f_\mathrm{slip}}_{\substack{\text{syringe-wall}\\\text{slip}}}
\cdot
\underbrace{R_\mathrm{rec}^{1/2}}_{\substack{\text{viscoelastic}\\\text{recovery}}}
\label{eq:kflow}
\end{equation}
where $\beta_\mathrm{swell}$ is the die-swell ratio (line width over nozzle ID,
minus one), $f_\mathrm{slip}$ is an empirical syringe-piston slip factor (default
1.0; calibrated experimentally), and $R_\mathrm{rec}$ is the structural recovery
ratio measured by \texttt{Recovery\_v1.py}. The Cross-model exponent $m$ is a
useful proxy for $\beta_\mathrm{swell}$: $m \to 1$ indicates a near-cooperative
shear-thinning collapse with low elastic energy storage and consequently low die
swell.

\subsection{Hagen--Poiseuille pressure for fitted models}
The extrusion pressure drop is integrated along the needle as
\begin{equation}
\Delta P \;=\; 4 L \frac{\tau_w}{D}, \qquad
\tau_w = \eta(\dot\gamma_w)\,\dot\gamma_w,
\label{eq:HP}
\end{equation}
with $\eta(\dot\gamma)$ given by the fitted Cross model
\begin{equation}
\eta(\dot\gamma) \;=\; \eta_\infty + \frac{\eta_0-\eta_\infty}{1+(K\,\dot\gamma)^{m}}.
\label{eq:cross}
\end{equation}
The values currently in use, summarised from
\texttt{rheology\_results}, are reproduced in Table~\ref{tab:rheo_params}.

\begin{table}[h]
\centering
\small
\caption{Cross-model parameters used by the MATLAB solver (Ramp~1, no pre-shear).}
\label{tab:rheo_params}
\begin{tabular}{l c c c c c}
\toprule
Ink & $\eta_0$ [Pa\,s] & $\eta_\infty$ [Pa\,s] & $K$ [s] & $m$ [-] & $R^2$ \\
\midrule
C15            & 578.3   & 0.0163 & 3.711  & 0.811 & 0.9956 \\
C15 Gira 5.5   & 2439.5  & 0.001  & 4.099  & 0.988 & 0.9882 \\
Bozano         & 6354.6  & 0.932  & 73.05  & 0.866 & 0.9937 \\
\bottomrule
\end{tabular}
\end{table}

% ====================================================================
\section{(a)~MATLAB post-processing: \texorpdfstring{$Q(\Delta P)$}{Q(dP)} to slicer lookup}
\label{sec:matlab}

The existing solver \texttt{bioprinting\_algorithm\_cross\_v2.m} writes a
$Q(\Delta P)$ array per sample. The block below post-processes that array into a
slicer-ready lookup table linking $\Delta P$, $\bar v_\mathrm{nozzle}$,
$v_\mathrm{print}$ (under the assumption $w=\mathrm{ID}$, $h=0.7\cdot\mathrm{ID}$),
and a predicted $k_\mathrm{flow}$ derived from Eq.~\ref{eq:kflow}.

\begin{lstlisting}[style=matlab,caption={Post-processing block to append to
\texttt{run\_solver\_Cross\_v2.m}. Produces a CSV
\texttt{slicer\_lookup\_<sample>.csv} per ink.},label={lst:matlab_postproc}]
% ============================================================
%  Post-processing: Q(dP)  -->  slicer-ready parameter table
%  Append after the existing solver loop in run_solver_Cross_v2.m
% ============================================================
%  Required inputs (already produced by the Cross solver):
%      dP_vec    : column of pressures used (Pa)
%      Q_vec     : column of corresponding volumetric flows (m^3/s)
%      ID        : nozzle inner diameter (m), e.g. 0.515e-3 for 21G
%      m_cross   : Cross exponent for this sample (dimensionless)
%      Rrec_pct  : structural recovery percentage at the deposition
%                  shear rate, from Recovery_v1.py
% ------------------------------------------------------------

w_line  = ID;                     % assumed 1:1 line-to-ID
h_layer = 0.7 * ID;               % typical DIW default (m)

% Mean nozzle velocity from continuity (m/s)
A_nozzle = pi * (ID/2)^2;
v_mean   = Q_vec ./ A_nozzle;

% Head speed for w*h deposition (m/s -> mm/s)
v_print  = (Q_vec ./ (w_line * h_layer)) * 1e3;

% Die-swell estimate from Cross exponent
%   beta -> 0 for cooperative shear-thinning collapse (m near 1)
%   beta -> 0.3 for mild thinning (m ~ 0.6)
beta_swell = max(0, 0.30 * (1 - m_cross));

% Flow factor (Eq. 3): k_flow = (1+beta)^2 * f_slip * sqrt(Rrec)
f_slip    = 1.0;                  % default; calibrate per syringe
k_flow    = (1 + beta_swell).^2 .* f_slip .* sqrt(Rrec_pct/100);

% Assemble table
T = table(dP_vec/1e3, v_mean*1e3, v_print, ...
          repmat(beta_swell, size(dP_vec)), ...
          repmat(k_flow,    size(dP_vec)), ...
    'VariableNames', {'dP_kPa','v_mean_mm_s','v_print_mm_s', ...
                      'beta_swell','k_flow'});

outname = sprintf('slicer_lookup_%s.csv', sample_name);
writetable(T, fullfile(output_folder, outname));
fprintf('[Saved] %s\n', outname);
\end{lstlisting}

\noindent\textbf{Worked example (C15~Gira~5.5, 21G needle, $\Delta P=252\,\mathrm{kPa}$).}
With $\mathrm{ID}=0.515\,\mathrm{mm}$, $h_\mathrm{layer}=0.40\,\mathrm{mm}$,
$m=0.988$, and recovery $R_\mathrm{rec}=31.2\%$ at $100\,\mathrm{Hz}$:
$\beta_\mathrm{swell}\approx 0.30(1-0.988)=0.004$;
$k_\mathrm{flow}\approx (1.004)^2\cdot 1\cdot\sqrt{0.312}=0.56$.
At the simulated mean velocity of $10\mms$,
$v_\mathrm{print}\approx 10\mms$ for $1{:}1$ deposition.

\noindent\textbf{Interpretation.} A $k_\mathrm{flow}$ value below unity means the
slicer's nominal volumetric command will \emph{under-deposit}; the firmware E-steps
must be increased by $1/0.56 \approx 1.78\times$, or equivalently the slicer's
``extrusion multiplier'' set to $\sim 1.78$. This is the rheologically grounded
starting point for the cube calibration in Section~\ref{sec:calibration}.

% ====================================================================
\section{(b)~Calibration-cube protocol (Materials \& Methods prose)}
\label{sec:calibration}

The following four-test sequence calibrates the FFF-style firmware to the
non-Newtonian, recoverable ink. Each test maps to a single slicer parameter so
that, if a print fails, the failure mode points directly to a single rheological
input.

\paragraph{Test C1 --- Print-speed sweep (single-line deposition).}
A $30\,\mathrm{mm}$ straight filament is deposited at five head speeds
($v_\mathrm{print}\in\{5,10,20,30,50\}\mms$) using the same commanded pressure
$\Delta P$ from the MATLAB lookup (Table~\ref{lst:matlab_postproc}). The deposited
filament diameter $d_\mathrm{fil}$ is measured with a calibrated stereoscope at
three points along each line and reported as
$d_\mathrm{fil}/\mathrm{ID}$. The optimal $v_\mathrm{print}$ is the value at which
$d_\mathrm{fil}/\mathrm{ID}$ lies between $0.95$ and $1.10$; lower values indicate
under-extrusion (head outrunning ink supply), higher values indicate
over-deposition or accumulated swell. The line width $w_\mathrm{line}$ used by the
slicer is set to $\mathrm{ID}\,(1+\beta_\mathrm{swell})$ with $\beta_\mathrm{swell}$
inferred from $d_\mathrm{fil}/\mathrm{ID}$ at the optimal speed.

\paragraph{Test C2 --- Mass-per-line calibration ($k_\mathrm{flow}$).}
At the optimal $v_\mathrm{print}$ identified in C1, ten parallel $30\,\mathrm{mm}$
lines are printed on a tared substrate. The substrate is reweighed; deposited mass
is converted to deposited volume using the ink density measured by displacement
($\rho_\mathrm{ink}=1.00\pm 0.02\,\mathrm{g/cm^3}$ for hydrogel inks). The
calibrated flow factor is
\begin{equation}
k_\mathrm{flow}^\mathrm{exp}
\;=\;
\frac{V_\mathrm{deposited}^\mathrm{measured}}
     {v_\mathrm{print}\cdot w_\mathrm{line}\cdot h_\mathrm{layer}\cdot t_\mathrm{print}}.
\end{equation}
This is the value entered into the slicer; it supersedes the Eq.~\ref{eq:kflow}
prediction once measured.

\paragraph{Test C3 --- Layer-height sweep (3-layer stack).}
A $3$-layer $10\times 10\,\mathrm{mm}$ square is printed at three layer heights
$h_\mathrm{layer}\in\{0.3,\,0.4,\,0.5\}\,\mathrm{mm}$ (i.e.\ $0.58$, $0.78$, $0.97$
of ID for 21G). After $30\,\mathrm{min}$ of relaxation, the printed stack height
is measured. The first $h_\mathrm{layer}$ at which the measured total height drops
below $90\%$ of the commanded total height defines the practical upper bound
$h_\mathrm{layer}^\mathrm{max}$. This bound is independently constrained by the
SAOS-derived predictions in Section~\ref{sec:saos_hmax} (Table~\ref{tab:hmax}).

\paragraph{Test C4 --- Vertical-tower test (scaffold-height bound).}
A $5\times 5\,\mathrm{mm}$ tower is printed at $h_\mathrm{layer}^\mathrm{max}$
from C3, layer by layer, up to $N=40$ layers. The tower is photographed every
$5$ layers and again after $30\,\mathrm{min}$. The number of layers at which the
base width increases by $>10\%$ defines the experimental maximum scaffold height
$h_\mathrm{scaff}^\mathrm{exp} = N_\mathrm{max}\cdot h_\mathrm{layer}^\mathrm{max}$,
which is then compared against the three SAOS-derived predictions.

\paragraph{Bridging tests for unsupported spans.}
For lattice scaffolds with unsupported struts, a bridging test prints horizontal
filaments over gaps of $L\in\{0.5,1.0,2.0\}\,\mathrm{mm}$. Maximum sag is
recorded; the buckling-corrected criterion (Eq.~\ref{eq:smay}) provides the
prediction.

\noindent\textbf{Reporting.} For each ink, a single table summarises the four
outcomes: the optimal $v_\mathrm{print}$, the measured $k_\mathrm{flow}$, the
calibrated $h_\mathrm{layer}^\mathrm{max}$, and the achieved scaffold height
$h_\mathrm{scaff}^\mathrm{exp}$ vs.\ the three SAOS predictions.

% ====================================================================
\section{(c)~SAOS-derived maximum scaffold height}
\label{sec:saos_hmax}

\subsection{Rationale and yield criteria}
The vertical fidelity of a DIW print is bounded by the requirement that the
bottom layer must support the gravitational stress of all subsequent layers
without yielding. The stacking stress is
\begin{equation}
\sigma_\mathrm{stack} \;=\; \rho\,g\,N\,h_\mathrm{layer},
\label{eq:stack_stress}
\end{equation}
with $\rho=1000\,\mathrm{kg/m^3}$ (hydrogel approximation),
$g=9.81\,\mathrm{m/s^2}$. The corresponding maximum stacked height
$h_\mathrm{max}=\sigma_\mathrm{max}/(\rho g)$ depends on which yield criterion is
adopted for $\sigma_\mathrm{max}$. We report three:
\begin{description}[leftmargin=2.4em]
\item[(A) LVR-endpoint criterion] $\sigma_\mathrm{max} =
G'_\mathrm{LVR}\cdot\gamma_\mathrm{LVR}$. The strictest interpretation: the bottom
layer must remain within the linear-viscoelastic regime.
\item[(B) Practical $G'(\omega=1\,\mathrm{rad/s})$ criterion]
$\sigma_\mathrm{max} = G'(\omega=1)$. Standard literature default
\citep{PaxtonEtAl2017, SchwabEtAl2020}: uses the storage modulus directly as a
yield-stress proxy at the print-relaxation timescale.
\item[(C) Quasi-static $G'(\omega=0.01\,\mathrm{rad/s})$ criterion]
$\sigma_\mathrm{max} = G'(\omega=0.01)$, obtained by low-frequency power-law
extrapolation of the SAOS sweep. This is the appropriate bound when the
inter-pass dwell exceeds $\sim 100\,\mathrm{s}$.
\end{description}

For unsupported spans (lattice struts), the Smay buckling criterion
\citep{HabibKhoda2018} gives
\begin{equation}
h_\mathrm{max}(L) \;=\; \sqrt{\frac{1.94\,G'}{\rho g}}\cdot\sqrt{L},
\label{eq:smay}
\end{equation}
which is tighter than Eq.~\ref{eq:stack_stress} for any $L\lesssim 5\,\mathrm{mm}$.

\subsection{Low-frequency extrapolation of \texorpdfstring{$G'(\omega)$}{G'(omega)}}
The frequency sweeps cover $\omega\in[1, 240]\,\mathrm{rad/s}$ (Anton Paar
AresG2, cone--plate 50\,mm, $0.1\,\mathrm{mm}$ gap, $25\,^\circ\mathrm{C}$). The
six lowest-frequency points were fitted in log space to
$G'(\omega)=G'_0\,\omega^{\beta}$ (Table~\ref{tab:lowomega_fit}). The exponent
$\beta$ encodes the gel character: $\beta\to 0$ for a fully formed gel (Bozano);
$\beta\to 0.5$ for a critical-gel polymer solution
\citep{WinterChambon1987} (C15 and the NE).

\begin{table}[h]
\centering
\small
\caption{Low-frequency power-law fit $G'(\omega) = G'_0\,\omega^{\beta}$ over the
six lowest measured points.}
\label{tab:lowomega_fit}
\begin{tabular}{l c c c c c}
\toprule
Ink & $G'_0$ [Pa] & $\beta$ [-] & $R^2$ & $G'(\omega=1)$ [Pa] & $G'(\omega=0.01)$ [Pa] \\
\midrule
C15            & 114.6 & 0.454 & 1.000 & 114.5 & 14.1 \\
C15 Gira 5.5   & 234.9 & 0.446 & 0.760 & 228.4 & 30.2 \\
Bozano         & 348.9 & 0.091 & 0.979 & 345.8 & 229.1 \\
\bottomrule
\end{tabular}
\end{table}

\noindent\textbf{Physical reading.} The exponents reveal a fundamental
microstructural distinction. Bozano is a true gel: $\beta=0.091$ implies a flat,
time-invariant elastic network. C15 and the NE inclusion sit closer to the
Winter--Chambon critical gel exponent ($\beta\approx 0.5$), consistent with
\emph{transient, physically associated} networks. Crucially, the NE addition
\emph{preserves} the critical-gel character ($\beta=0.446$ vs.~$0.454$ for the
pure CMC base) while raising the absolute modulus by $\sim 2\times$
($G'_0=234.9\Pa$ vs.~$114.6\Pa$). The NE acts as a \emph{filler}, not as a
crosslinker --- a key positioning argument for the manuscript.

\subsection{Predicted maximum scaffold height per ink}

\begin{table}[h]
\centering
\small
\caption{Maximum scaffold height $h_\mathrm{max}$ (stacked, full sub-layer
support) under three SAOS-derived yield criteria, and the buckling-corrected
unsupported-span bound for $L=1\,\mathrm{mm}$ (Eq.~\ref{eq:smay}).
$\rho=1000\,\mathrm{kg/m^3}$, $g=9.81\,\mathrm{m/s^2}$.}
\label{tab:hmax}
\begin{tabular}{l c c c c c c c}
\toprule
\multirow{2}{*}{Ink}
 & \multicolumn{2}{c}{(A) LVR endpoint}
 & \multicolumn{2}{c}{(B) $G'(\omega=1)$}
 & \multicolumn{2}{c}{(C) $G'(\omega=0.01)$}
 & Buckle \\
\cmidrule(lr){2-3}\cmidrule(lr){4-5}\cmidrule(lr){6-7}\cmidrule(lr){8-8}
 & $\sigma$ [Pa] & $h$ [mm] & $\sigma$ [Pa] & $h$ [mm] & $\sigma$ [Pa] & $h$ [mm] & $h$(1\,mm)[mm] \\
\midrule
C15           & 172.9 & 17.6 & 114.5 & 11.7 & 14.1  & 1.4  & 4.76 \\
C15 Gira 5.5  & \textbf{558.5} & \textbf{56.9} & 228.4 & 23.3 & 30.2  & 3.1  & 6.72 \\
Bozano        & 47.6  & 4.8  & 345.8 & 35.3 & 229.1 & 23.4 & 8.27 \\
\bottomrule
\end{tabular}
\end{table}

\noindent\textbf{Three findings to highlight in the manuscript.}
\begin{enumerate}[leftmargin=1.6em]
\item \emph{The NE inclusion increases scaffold height under every short-timescale
criterion.} On the LVR-endpoint bound, the NE raises $h_\mathrm{max}$ from
$17.6\,\mathrm{mm}$ to $56.9\,\mathrm{mm}$ ($3.2\times$). This is because the NE
inclusion simultaneously raises $G'$ ($114\to 228\Pa$) \emph{and} preserves the
ink's broad LVR ($\gamma_\mathrm{LVR}=61\%$ vs.\ $64\%$ for pure C15).
\item \emph{Bozano is a slow-timescale outlier.} The narrow Bozano LVR
($\gamma_\mathrm{LVR}=13.7\%$) gives the smallest LVR-endpoint $h_\mathrm{max}$
($4.8\,\mathrm{mm}$), but its true-gel character ($\beta=0.09$) means it retains
$66\%$ of its $\omega=1\,\mathrm{rad/s}$ modulus at $\omega=0.01\,\mathrm{rad/s}$,
delivering the highest quasi-static $h_\mathrm{max}$ ($23.4\,\mathrm{mm}$). For
prints with long dwell times the Bozano network is preferred; for short-dwell
multi-layer prints the NE wins.
\item \emph{Buckling, not stacking, is the binding constraint for lattices.}
For $1\,\mathrm{mm}$ unsupported spans (Eq.~\ref{eq:smay}), all three inks lie in
the $5$--$8\,\mathrm{mm}$ band; even the NE drops from $57\,\mathrm{mm}$
(stacked) to $6.7\,\mathrm{mm}$ (buckling-limited). Lattice scaffold designs
must therefore use shorter struts or higher modulus inks (e.g.\ blended
formulations).
\end{enumerate}

% ====================================================================
\section{Synthesis: predicted slicer parameters per ink}
\label{sec:synthesis}

Table~\ref{tab:slicer_predicted} compiles, for each ink, the recommended starting
slicer parameters --- to be refined by the Section~\ref{sec:calibration}
calibration cubes during the printing-validation campaign.

\begin{table}[h]
\centering
\small
\caption{Predicted DIW slicer parameters for a 21G needle (ID = 0.515\,mm,
$L = 31.75\,\mathrm{mm}$) at the simulated mean nozzle velocity
$\bar v_\mathrm{nozzle}=10\mms$.}
\label{tab:slicer_predicted}
\begin{tabular}{l c c c}
\toprule
Parameter & C15 & C15 Gira 5.5 (NE) & Bozano \\
\midrule
$\Delta P_\mathrm{ext}$ [kPa]            & 1041   & \textbf{252}    & --       \\
Wall shear $\tau_w$ [Pa]                  & 4197   & \textbf{963}    & --       \\
$v_\mathrm{print}$ [mm/s]                 & 10     & 10              & 10        \\
$h_\mathrm{layer}$ [mm]                   & 0.35   & 0.40            & 0.40     \\
$w_\mathrm{line}$ [mm]                    & 0.52   & 0.52            & 0.57     \\
$k_\mathrm{flow}$ (predicted)             & 0.82   & 0.56            & 1.09     \\
$v_\mathrm{travel}$ [mm/s]                & 50--80 & 50--80          & 80--120  \\
$h_\mathrm{scaffold}$ (stack, LVR) [mm]   & 17.6   & \textbf{56.9}   & 4.8      \\
$h_\mathrm{scaffold}$ (stack, $\omega=1$) [mm] & 11.7 & 23.3 & 35.3 \\
$h_\mathrm{scaffold}$ (buckle, $L=1$\,mm) [mm] & 4.8 & 6.7 & 8.3 \\
\bottomrule
\end{tabular}
\end{table}

\paragraph{Bottom-line guidance.}
For the present manuscript scope (CMC + sunflower-oil NE, 21G DIW), the entries
in bold define the central narrative: the NE formulation prints with
$\sim\!1/4$ of the extrusion pressure of pure C15, with cell-safe wall shear
($<5\,\mathrm{kPa}$), and supports a $\sim\!3\times$ taller stacked scaffold
under the LVR criterion. The buckling and quasi-static bounds remain similar to
the C15 baseline, identifying the next experimental priority: validate the LVR
prediction with a calibrated vertical-tower print.

% ====================================================================
\section{Open items}
\label{sec:open}

Three measurements would close the framework:
\begin{enumerate}[leftmargin=1.6em]
\item \emph{Single direct die-swell measurement} (extrude into air, measure
$d_\mathrm{fil}/\mathrm{ID}$) per ink, to replace the Cross-exponent $\beta_\mathrm{swell}$
proxy in Eq.~\ref{eq:kflow}.
\item \emph{Ink density by displacement} per ink, replacing the
$1.00\,\mathrm{g/cm^3}$ assumption in $k_\mathrm{flow}^\mathrm{exp}$.
\item \emph{Direct yield-stress measurement} (Herschel--Bulkley refit of all flow
data, or stress-controlled creep test), to corroborate the SAOS-derived
$\sigma_\mathrm{max}$ values in Table~\ref{tab:hmax}.
\end{enumerate}

% ====================================================================
\section*{Script and data provenance}

\begin{itemize}[leftmargin=1.6em]
\item \texttt{Fit\_Muitos\_Modelos\_v4.py} (\emph{active}): flow-curve fitting;
   produced Cross parameters of Table~\ref{tab:rheo_params}.
\item \texttt{Recovery\_v1.py} (\emph{active}): structural recovery used in
   Eq.~\ref{eq:kflow}; output in
   \texttt{Analises/Python/Results/Relatorio\_Recovery.txt}.
\item \texttt{extract\_SAOS\_values.py} (\emph{active}): amplitude- and
   frequency-sweep parser; $G'_\mathrm{LVR}$ and $\gamma_\mathrm{LVR}$ source for
   Table~\ref{tab:hmax}.
\item \texttt{extract\_hmax\_v1.py} (\emph{new, v1, active}): low-$\omega$
   extrapolation and $h_\mathrm{max}$ computation; produced
   Tables~\ref{tab:lowomega_fit} and~\ref{tab:hmax}; output in
   \texttt{Analises/Python/Results/SAOS\_hmax\_v1.txt}.
\item \texttt{bioprinting\_algorithm\_cross\_v2.m} +
   \texttt{run\_solver\_Cross\_v2.m} (\emph{active}): Cross-model extrusion
   solver; receives the MATLAB block of Listing~\ref{lst:matlab_postproc}.
\end{itemize}

% ====================================================================
\bibliographystyle{abbrvnat}
\bibliography{references}

\end{document}
